\section{Engel Theorem the third type}

\begin{definition}
    Let $L$ be a Lie algebra over a field $F$. An $L$-module is a vector space $V$ over $F$ together with a bilinear map
    \[
        L\times V \to V,\qquad (x,v)\mapsto x\cdot v,
    \]
    such that
    \[
        [x,y]\cdot v = x\cdot (y\cdot v)-y\cdot (x\cdot v)
    \]
    for all $x,y\in L$ and $v\in V$. In other words, $V$ is an $L$-module if and only if there exists a Lie algebra homomorphism
    \[
        \phi:L\to \mathfrak{gl}(V).
    \]
    such that
    \[
        \phi(x)(v) = x\cdot v
    \]
    for all $x\in L$ and $v\in V$.
\end{definition}

\begin{proposition}
    Let $L$ be a Lie algebra over a field $F$, and let $M$ be a vector space over $F$. Suppose a bilinear operation
    \[
        L\times M\to M,\qquad (a,m)\mapsto a\cdot m
    \]
    is given. Then the following conditions are equivalent:

    \begin{enumerate}
        \item The vector space $L\oplus M$ becomes a Lie algebra under the bracket defined by
              \[
                  [a+m,\ b+n]=[a,b]+a\cdot n-b\cdot m
                  \qquad (a,b\in L,\ m,n\in M),
              \]
              equivalently,
              \[
                  [a,m]=a\cdot m,\qquad [m,a]=-\,a\cdot m
              \]
              for all $a\in L$, $m\in M$, and
              \[
                  [m,n]=0
              \]
              for all $m,n\in M$.

        \item For all $a,b\in L$ and $m\in M$,
              \[
                  [a,b]\cdot m=a\cdot (b\cdot m)-b\cdot (a\cdot m).
              \]

        \item The map
              \[
                  \rho:L\to \mathfrak{gl}_F(M),\qquad \rho(a)(m)=a\cdot m,
              \]
              is a Lie algebra homomorphism.
    \end{enumerate}

    When these equivalent conditions hold, we say that $M$ is an $L$-module.
\end{proposition}

\begin{proof}
    We prove $(1)\Rightarrow(2)\Rightarrow(3)\Rightarrow(1)$.

    \medskip

    $(1)\Rightarrow(2)$. Assume that $L\oplus M$ is a Lie algebra with the stated bracket. Since every Lie algebra satisfies the Jacobi identity, for any $a,b\in L$ and $m\in M$ we have
    \[
        [a,[b,m]]+[b,[m,a]]+[m,[a,b]]=0.
    \]
    Now
    \[
        [b,m]=b\cdot m,\qquad [m,a]=-\,a\cdot m,\qquad [a,b]\in L.
    \]
    Because $M$ is abelian, this gives
    \[
        [a,b\cdot m]+[b,-\,a\cdot m]+[m,[a,b]]=0.
    \]
    Using again the definition of the bracket,
    \[
        a\cdot (b\cdot m)-b\cdot (a\cdot m)-[a,b]\cdot m=0.
    \]
    Hence
    \[
        [a,b]\cdot m=a\cdot (b\cdot m)-b\cdot (a\cdot m).
    \]

    \medskip

    $(2)\Rightarrow(3)$. Define
    \[
        \rho:L\to \mathfrak{gl}_F(M),\qquad \rho(a)(m)=a\cdot m.
    \]
    Since the action is bilinear, $\rho$ is linear. For any $a,b\in L$ and $m\in M$, condition (2) gives
    \[
        \rho([a,b])(m)=[a,b]\cdot m
        =a\cdot (b\cdot m)-b\cdot (a\cdot m)
        =\rho(a)\rho(b)(m)-\rho(b)\rho(a)(m).
    \]
    Therefore
    \[
        \rho([a,b])=[\rho(a),\rho(b)],
    \]
    so $\rho$ is a Lie algebra homomorphism.

    \medskip

    $(3)\Rightarrow(1)$. Assume that $\rho:L\to \mathfrak{gl}_F(M)$ is a Lie algebra homomorphism, and define a bracket on $L\oplus M$ by
    \[
        [a+m,\ b+n]=[a,b]+\rho(a)(n)-\rho(b)(m).
    \]
    This bracket is bilinear and skew-symmetric. It remains to verify the Jacobi identity.

    Take $a,b,c\in L$ and $m,n,p\in M$. Write
    \[
        x=a+m,\qquad y=b+n,\qquad z=c+p.
    \]
    A direct calculation shows that the $L$-component of
    \[
        [x,[y,z]]+[y,[z,x]]+[z,[x,y]]
    \]
    is
    \[
        [a,[b,c]]+[b,[c,a]]+[c,[a,b]],
    \]
    which is zero because $L$ is a Lie algebra.

    For the $M$-component, expanding the definition of the bracket yields
    \[
        \rho(a)\rho(b)(p)-\rho(a)\rho(c)(n)-\rho([b,c])(m)
        +\rho(b)\rho(c)(m)-\rho(b)\rho(a)(p)-\rho([c,a])(n)
    \]
    \[
        +\rho(c)\rho(a)(n)-\rho(c)\rho(b)(m)-\rho([a,b])(p).
    \]
    Grouping the terms by $m,n,p$, we obtain
    \[
        \bigl(\rho(b)\rho(c)-\rho(c)\rho(b)-\rho([b,c])\bigr)(m)
    \]
    \[
        +\bigl(\rho(c)\rho(a)-\rho(a)\rho(c)-\rho([c,a])\bigr)(n)
    \]
    \[
        +\bigl(\rho(a)\rho(b)-\rho(b)\rho(a)-\rho([a,b])\bigr)(p).
    \]
    Since $\rho$ is a Lie algebra homomorphism, each bracketed expression is zero. Hence the $M$-component also vanishes, so the Jacobi identity holds.

    Therefore $L\oplus M$ is a Lie algebra with the stated bracket.

    This proves that the three conditions are equivalent.
\end{proof}

\begin{theorem}
    Let $L$ be a finite dimensional Lie algebra over a field $F$. If $\forall a \in L$, $\operatorname{ad}(a) $ is nilpotent, then $L$ is nilpotent.
\end{theorem}

\begin{proof}
    Note that since $L$ is finite-dimensional, say $\dim L=n<\infty$, and for every $a\in L$ we have
    \[
        (\operatorname{ad} a)^n=0.
    \]
    Hence we may apply the previous theorem: if for all $a\in L$,
    \[
        (\operatorname{ad} a)^n=0,
    \]
    then
    \[
        (\operatorname{ad} L)^{\,n}=0,
    \]
    and therefore
    \[
        L^{n+1}=0.
    \]
    So $L$ is nilpotent.
\end{proof}

\begin{remark}
    If $L$ is infinite-dimensional, one should distinguish two cases:
    \[
        \left\{
        \begin{array}{l}
            \operatorname{ad} a \text{ nilpotent with a uniform bound } n
            \ \Longrightarrow\
            L \text{ nilpotent}, \\[4pt]
            \text{no uniform bound on the nilpotency indices}
            \ \nRightarrow\
            L \text{ nilpotent}.
        \end{array}
        \right.
    \]
\end{remark}

\begin{remark}
    There are three related versions, with
    \[
        \text{type }(2)\Longrightarrow \text{type }(1)\Longrightarrow \text{type }(3).
    \]

    Type (2): Let $S\subseteq L$ be a Lie subset. And $S$ generates $L$ as a Lie algebra. Assume there exists $n$ such that for every $a\in S$
    \[
        a^n =0.
    \]
    Then $A=<L>$ is a nilpotent algebra. (When $A$ is a finite-dimensional associative algebra).

    In particular, if we choose $S$ to be $L$, then we get type (1).

    Type (1): Let $ A $ be a finite dimensional enveloping algebra of $L$, i.e. $ L\subseteq A^{(-)} $, $ A = <L> $.If $ \exists n \ge 1 $ such that $ \forall a \in L $, $ a^n=0 $. Then $ A $ is nilpotent.

    And consider $A$ as $\operatorname{ad}(L)$, we get type (3).

    Type (3):    Let $L$ be a finite dimensional Lie algebra over a field $F$. If $\forall a \in L$, $\operatorname{ad}(a) $ is nilpotent, then $L$ is nilpotent.
\end{remark}

We also have the following truth:

\begin{corollary}
    \ Let $M$ be an $L$-module with $\dim M<\infty$, that is, let
    \[
        \varphi:L\to \mathfrak{gl}_F(M)
    \]
    be a representation. Assume that for every $a\in L$, the operator $\varphi(a)$ is nilpotent. Equivalently, for every $a\in L$ there exists $n$ such that
    \[
        \underbrace{a\bigl(a(\cdots(a}_{n\text{ times}}\,M)\cdots)\bigr)=0.
    \]
    Then there exists $N$ such that
    \[
        \underbrace{L\cdots L}_{N\text{ times}}M=0.
    \]
    In other words, if every $\varphi(a)$ is nilpotent, then the Lie algebra $\varphi(L)$ is nilpotent. Hence there exists a basis of $M$ in which every element of $\varphi(L)$ is represented by a strictly upper triangular matrix:
    \[
        \varphi(L)\subseteq
        \left\{
        \begin{pmatrix}
            0      & *      & \cdots & *      \\
            0      & 0      & \cdots & *      \\
            \vdots & \vdots & \ddots & \vdots \\
            0      & 0      & \cdots & 0
        \end{pmatrix}
        \right\}.
    \]
\end{corollary}

\begin{remark}
    This corollary is the module-theoretic form of Engel's theorem. Indeed, the representation
    \[
        \varphi:L\to \mathfrak{gl}_F(M)
    \]
    has image $\varphi(L)$, which is a Lie subalgebra of $\mathfrak{gl}_F(M)$. By assumption, every element of $\varphi(L)$ is a nilpotent endomorphism of the finite-dimensional vector space $M$. Hence, by Engel's theorem, the Lie algebra $\varphi(L)$ is nilpotent.

    The conclusion
    \[
        \underbrace{L\cdots L}_{N\text{ times}}M=0
    \]
    means that there exists a uniform integer $N$ such that for all $a_1,\dots,a_N\in L$ and all $m\in M$,
    \[
        a_1\bigl(a_2(\cdots(a_Nm)\cdots)\bigr)=0.
    \]
    Thus the action is nilpotent not only for each fixed $a\in L$, but uniformly for all sufficiently long iterated actions.

    A further consequence of Engel's theorem is that there exists a basis of $M$ with respect to which every operator $\varphi(a)$ is represented by a strictly upper triangular matrix. Equivalently, $M$ admits a chain of subspaces
    \[
        0=M_0\subset M_1\subset \cdots \subset M_r=M,
        \qquad \dim M_i=i,
    \]
    such that
    \[
        L\cdot M_i\subseteq M_{i-1}
        \qquad (1\le i\le r).
    \]

    Explanation: By Engel's theorem, $\varphi(L)$ is nilpotent. Hence
    \[
        \underbrace{L\cdots L}_{N\text{ times}}M=0
    \]
    Take the minimal $N$ such that $L^N M=0$ and $L^{N-1}M \neq 0$. Then we have a non-zero vector $v_1$ in $L^{N-1}M$ which is killed by $\varphi(L)$. Denote $M_1 = \operatorname{span}(v_1)$. Then $L\cdot M_1 \subseteq M_0$. Now consider $M/M_1$. By Engel's theorem and by induction hypothesis, $\varphi(L)$ is nilpotent on $M/M_1$. Hence there exists a basis of $M/M_1$ with respect to which every element of $\varphi(L)$ is represented by a strictly upper triangular matrix. Equivalently, $M/M_1$ admits a chain of subspaces
    \[
        0=M_1/M_1\subset M_2/M_1\subset \cdots \subset M_r/M_1=M/M_1,
        \qquad \dim M_i/M_1=i,
    \]
    such that
    \[
        L\cdot (M_i/M_1)\subseteq M_{i-1}/M_1
        \qquad (1\le i\le r).
    \]
    Now we can take a basis of $M_1$ and $M/M_1$ to get a basis of $M$ with respect to which every element of $\varphi(L)$ is represented by a strictly upper triangular matrix.

    In particular, this corollary reduces to the usual form of Engel's theorem when one takes $M=L$ and $\varphi=\operatorname{ad}$.
\end{remark}

\begin{remark}
    According to this corollary and type (2), we can rewrite it as:

    Let $M$ be an $L$-module with $\dim M<\infty$, that is, let
    \[
        \varphi:L\to \mathfrak{gl}_F(M)
    \]
    be a representation. Assume that for every $a\in S$, where $S$ is a Lie subset of $L$ and $<S>=L$, the operator $\varphi(a)$ is nilpotent. Equivalently, for every $a\in S$ there exists $n$ such that
    \[
        \underbrace{a\bigl(a(\cdots(a}_{n\text{ times}}\,M)\cdots)\bigr)=0.
    \]
    Then there exists $N$ such that
    \[
        \underbrace{L\cdots L}_{N\text{ times}}M=0.
    \]
    For this result, we do not need $L$ or the enveloping algebra $A=<L>$ to be finite-dimensional now, we only need such $n$ is a universal bound.
\end{remark}

